\documentclass[12pt]{article}
\usepackage[utf8]{inputenc}
\usepackage[spanish]{babel}
\usepackage{amsmath}
\usepackage{amssymb}
\usepackage{booktabs}
\usepackage{geometry}
\geometry{margin=2.5cm}
\usepackage[most]{tcolorbox}
\newtcolorbox{intuicion}{
  colback=blue!5!white,
  colframe=blue!40!black,
  fonttitle=\bfseries,
  title=Intuici\'on,
  boxrule=0.6pt,
  arc=4pt,
  left=6pt, right=6pt, top=4pt, bottom=4pt
}
\newtcolorbox{intuicionmat}{
  colback=green!5!white,
  colframe=green!40!black,
  fonttitle=\bfseries,
  title=\textit{?`Qu\'e dice esta ecuaci\'on?},
  boxrule=0.6pt,
  arc=4pt,
  left=6pt, right=6pt, top=4pt, bottom=4pt
}
\newtcolorbox{explicacion}{
  colback=orange!5!white,
  colframe=orange!50!black,
  fonttitle=\bfseries,
  title=Explicaci\'on,
  boxrule=0.6pt,
  arc=4pt,
  left=6pt, right=6pt, top=4pt, bottom=4pt
}

\title{\textbf{Resumen: Distribuci\'on de los Ingresos Corporativos\\
entre Dividendos, Ganancias Retenidas e Impuestos}\\[4pt]
\large \textit{Distribution of Incomes of Corporations Among Dividends,}\\
\textit{Retained Earnings, and Taxes}\\[4pt]
\normalsize John Lintner, \textit{The American Economic Review}, Vol.~46, N.\textdegree~2 (1956), pp.~97--113}
\author{}
\date{Junio 2026}

\begin{document}
\maketitle
\tableofcontents
\newpage

%%% =========================================================
\section{Introducci\'on y Motivaci\'on del Estudio}
%%% =========================================================

El art\'iculo de Lintner (1956) presenta los resultados de un estudio emp\'irico sobre la
\textbf{pol\'itica de dividendos corporativos (corporate dividend policy)}, con implicaciones
directas para las \textbf{fluctuaciones c\'iclicas (cyclical fluctuations)} y las
\textbf{tendencias de crecimiento de largo plazo (long-term growth trends)} de la econom\'ia.

La investigaci\'on procede en dos etapas. Primero revisa los hallazgos del trabajo de campo
(\textit{field investigations}) que son m\'as relevantes para el tema. Luego usa esos hallazgos
para construir un \textbf{modelo te\'orico (theoretical model)} del comportamiento de dividendos
corporativos y lo somete a pruebas estad\'isticas.

Un resultado clave de partida: las ganancias retenidas en un per\'iodo son en gran medida
\textit{subproducto} de las decisiones de dividendos, no el resultado de decisiones aut\'onomas
de ahorro. De manera similar, el efecto de los impuestos sobre el volumen de ahorro neto
corporativo opera principalmente a trav\'es de su impacto sobre la magnitud de las ganancias
netas, que son el determinante primario del volumen de dividendos.

\begin{intuicion}
El punto de partida es sorprendente: la empresa no decide cu\'anto ahorrar; decide
cu\'anto pagar como dividendo. Lo que queda es el ahorro. En consecuencia, para entender
c\'omo las empresas distribuyen su ingreso, hay que enfocarse en c\'omo toman la decisi\'on
de dividendos, no en la de ahorro.
\end{intuicion}

%%% =========================================================
\section{Metodolog\'ia del Trabajo de Campo (Field Investigation)}
%%% =========================================================

\subsection{Selecci\'on de Empresas}

Se revisaron m\'as de 600 empresas listadas y consolidadas, y se seleccionaron
\textbf{28 empresas industriales} para investigaci\'on detallada, garantizando un m\'inimo
de 3 empresas por cada combinaci\'on de caracter\'isticas relevantes. Los criterios de
selecci\'on incluyeron:

\begin{itemize}
  \item Nivel de gasto en planta y equipo en la posguerra.
  \item \textbf{Raz\'on de pago (pay-out ratio)}: desde empresas que pagaron m\'as del 70\,\%
        de sus ganancias hasta aquellas que pagaron menos del 40\,\%.
  \item Uso de financiamiento externo durante el per\'iodo.
  \item Tama\~no de la empresa, frecuencia de cambios en las tasas, liquidez, estabilidad
        de ganancias, capitalizaci\'on y estructura de propiedad.
\end{itemize}

Las empresas seleccionadas pertenecen al sector \textbf{industrial (industrial)}, que exhibe
mayor diversidad en pol\'iticas de dividendos que otros sectores.

\begin{intuicion}
Las 28 empresas no son una muestra aleatoria para hacer inferencia estad\'istica cl\'asica.
Son una muestra \textit{dise\~nada} para capturar contrastes instructivos entre empresas
similares en algunos aspectos pero distintas en otros. Esto hace que los hallazgos sean
cualitativamente robustos, aunque no estad\'isticamente representativos por s\'i solos.
\end{intuicion}

\subsection{Recopilaci\'on de Informaci\'on}

Se realiz\'o un an\'alisis financiero completo con fuentes p\'ublicas y se realizaron
entrevistas con 2 a 5 directivos responsables por empresa (presidentes, vicepresidentes
financieros, tesoreros, contralores y directores). El an\'alisis cubri\'o 196
\textbf{a\~nos-empresa (company-years)} durante el per\'iodo 1947--1953.

\begin{intuicion}
El m\'etodo combina datos duros (estados financieros) con informaci\'on cualitativa
(entrevistas). Esto es clave porque Lintner quiere entender \textit{c\'omo piensan}
los gerentes cuando deciden sobre dividendos, algo que los balances solos no revelan.
\end{intuicion}

%%% =========================================================
\section{Patrones Observados en la Pol\'itica de Dividendos}
%%% =========================================================

\subsection{La Pregunta Central de las Decisiones de Dividendos}

El hallazgo m\'as general y robusto es que los gerentes \textbf{no} piensan en la pol\'itica
de dividendos como la elecci\'on del nivel \'optimo a pagar en cada per\'iodo. En los
196 a\~nos-empresa estudiados, no se encontr\'o ninguna instancia en que la pregunta
``\textit{?`cu\'anto pagar?}'' se considerara sin referencia a la tasa preexistente de dividendos.

La pregunta que los gerentes se hacen en la pr\'actica es:
\begin{center}
\textit{?`Debemos cambiar la tasa actual de dividendos? Y si s\'i, ?`en cu\'anto?}
\end{center}

La \textbf{variable dependiente (dependent variable)} real en el proceso de decisi\'on
es el \textbf{cambio en la tasa existente ($\Delta D$)}, no el nivel absoluto del nuevo pago.

\begin{intuicion}
Los modelos te\'oricos est\'aticos de la \'epoca supon\'ian que la empresa optimiza el
dividendo de cada per\'iodo de forma independiente. Lintner muestra que en la pr\'actica
eso no ocurre: la tasa anterior es siempre el punto de partida, y el gerente s\'olo decide
si la cambia y cu\'anto. Esto implica \textit{inercia} estructural en los dividendos.
\end{intuicion}

\subsection{Inercia y Conservadurismo (Inertia and Conservatism)}

Las empresas exhiben una fuerte preferencia por:
\begin{itemize}
  \item Tasas de dividendo \textbf{razonablemente estables (reasonably stable)}, o con
        crecimiento gradual.
  \item Evitar cambios que luego deban revertirse en el corto plazo.
  \item El mercado premia la estabilidad o el crecimiento gradual de las tasas de dividendo.
\end{itemize}

Este conservadurismo conduce al patr\'on de \textbf{adaptaci\'on parcial (partial adaptation)}:
en lugar de ajustar los dividendos completamente al cambio en ganancias en el mismo a\~no, las
empresas ajustan solo una \textit{fracci\'on} del cambio indicado, y completan el ajuste a lo
largo de los a\~nos siguientes.

\begin{intuicion}
Imagine que las ganancias de la empresa suben 30\,\%. Si la empresa sube los dividendos
completamente ese a\~no y al siguiente las ganancias bajan, deber\'a recortar dividendos,
lo que los accionistas detestan. Para evitar ese riesgo, la empresa sube los dividendos
solo parcialmente y espera confirmar que el aumento en ganancias es permanente.
\end{intuicion}

\subsection{Las Ganancias como Factor Determinante Principal}

Las ganancias netas corrientes (\textbf{current net earnings}) son invariablemente el punto
de partida de la deliberaci\'on gerencial sobre si cambiar los dividendos. Otros factores
---como restricciones de endeudamiento, posici\'on de liquidez, oportunidades de inversi\'on---
dominaron las decisiones de dividendos en menos del 5\,\% de los a\~nos-empresa estudiados.

Razones por las que las ganancias dominan:
\begin{enumerate}
  \item Son reportadas frecuentemente y reciben amplia publicidad financiera.
  \item Los directivos consideran que los accionistas tienen inter\'es propietario
        en las ganancias y esperan recibir una ``parte justa'' en dividendos.
  \item Las razones basadas en ganancias son f\'acilmente comprensibles y persuasivas
        para accionistas y la comunidad financiera.
  \item Una pol\'itica basada en otras consideraciones resultar\'ia dif\'icil de explicar
        y parecer\'ia errática (\textit{ad hoc}).
\end{enumerate}

\begin{intuicion}
Las ganancias funcionan como la \textit{se\~nal} que activa y orienta la decisi\'on de
dividendos. Sin un cambio significativo en ganancias, la gerencia simplemente no busca
ni considera otros datos para decidir sobre dividendos.
\end{intuicion}

\subsection{Raz\'on Objetivo de Pago y Velocidad de Ajuste (Target Pay-out Ratio and Speed of Adjustment)}

Aproximadamente dos tercios de las empresas estudiadas ten\'ian una
\textbf{raz\'on objetivo de pago ($r$) (target pay-out ratio)} relativamente definida:
la proporci\'on de ganancias que, como meta de largo plazo, deber\'ia distribuirse como
dividendo. Esta raz\'on no es un l\'imite r\'igido a\~no a a\~no, sino un \textit{objetivo
o ideal} hacia el cual la empresa tiende a moverse.

Las razones objetivo observadas variaron de un 20\,\% a un 80\,\%, siendo el 50\,\% la
cifra m\'as com\'un. La mayor\'ia de las dem\'as apuntaba al 40\,\% o al 60\,\%.

Adem\'as de la raz\'on objetivo, las empresas ten\'ian est\'andares sobre la
\textbf{velocidad de ajuste ($c$) (speed of adjustment)}: la fracci\'on de la brecha entre
el dividendo objetivo y el dividendo del per\'iodo anterior que ajustan en el a\~no corriente.
Dos empresas buscaban el ajuste completo en un a\~no; la mayor\'ia ajustaba solo una fracci\'on
(entre 1/6 y 1/2 por a\~no).

\begin{intuicion}
Piense en el dividendo como un nivel de agua en un tanque. La raz\'on objetivo dice cu\'al
es el nivel ideal. La velocidad de ajuste dice qu\'e tan r\'apido abre la v\'alvula para
llegar a ese nivel. Una velocidad baja significa que la empresa llegar\'a al nivel ideal
muy despacio, a lo largo de varios a\~nos.
\end{intuicion}

\subsection{Determinantes de la Raz\'on Objetivo y la Velocidad de Ajuste}

Los factores que moldearon estos par\'ametros en cada empresa incluyen:

\begin{itemize}
  \item Perspectivas de crecimiento de la industria y de la empresa.
  \item Movimiento c\'iclico promedio de las oportunidades de inversi\'on y requerimientos
        de capital de trabajo.
  \item Importancia relativa de las ganancias de capital (\textit{capital gains}) para
        los accionistas frente al ingreso corriente por dividendos.
  \item La preferencia de los accionistas entre tasas estables o fluctuantes.
  \item El premio que el mercado asigna a la estabilidad en la tasa de dividendo.
  \item Solidez financiera de la empresa y pol\'iticas respecto al uso de deuda y
        nuevas emisiones de acciones.
\end{itemize}

Una vez establecidos mediante experiencia acumulada, estos est\'andares permanecieron
pr\'acticamente \textbf{invariantes en el tiempo (invariant over time)} para la gran mayor\'ia
de las empresas durante todo el per\'iodo de posguerra.

\begin{intuicion}
La pol\'itica de dividendos no se renegoci\'o a\~no a a\~no: se fij\'o hist\'oricamente
y luego se aplic\'o de forma mec\'anica. Esto es clave para el modelo: si los par\'ametros
son estables, se pueden estimar de datos hist\'oricos y usar para predecir.
\end{intuicion}

\subsection{El Papel de los Impuestos y la Liquidez}

Los impuestos afectan los dividendos \textit{indirectamente}: reducen las ganancias netas
reportadas y, como los est\'andares de dividendo se fijan en t\'erminos de ganancias netas
(despu\'es de impuestos), una mayor carga tributaria reduce mec\'anicamente el nivel de
dividendos que resulta de aplicar la raz\'on objetivo.

La posici\'on de liquidez y las oportunidades de inversi\'on corrientes influyen en los
est\'andares de largo plazo (raz\'on objetivo y velocidad de ajuste) pero rara vez
intervienen directamente en las decisiones a\~no a a\~no. La gerencia planifica su
\textbf{presupuesto de capital (capital budget)} de modo que llevar a cabo la pol\'itica de
dividendos establecida no la deje en posiciones de iliquidez severa.

\begin{intuicion}
Los impuestos no hacen que los gerentes recorten dividendos de forma activa; simplemente
reducen el ingreso que sirve de base para calcular el dividendo objetivo. La liquidez tampoco
``invade'' la decisi\'on de dividendos directamente: la empresa se asegura de que su
presupuesto de inversi\'on sea compatible con mantener el dividendo comprometido.
\end{intuicion}

%%% =========================================================
\section{El Modelo Te\'orico de Comportamiento de Dividendos}
%%% =========================================================

\subsection{Ecuaci\'on Fundamental del Modelo (Fundamental Equation)}

Los patrones de toma de decisiones observados en el trabajo de campo se pueden sintetizar en
un modelo te\'orico simple sujeto a prueba estad\'istica. El modelo propone que el
\textbf{cambio en dividendos ($\Delta D_{it}$)} de la empresa $i$ en el per\'iodo $t$
obedece a:

\begin{equation}
  \Delta D_{it} \;=\; a_i + c_i\!\left(D^*_{it} - D_{i,t-1}\right) + u_{it}
\end{equation}

donde los componentes son:

\begin{itemize}
  \item $D^*_{it} = r_i P_{it}$: el \textbf{dividendo objetivo (target dividend)}, igual a la
        raz\'on objetivo $r_i$ aplicada a las ganancias corrientes $P_{it}$.
  \item $D_{i,t-1}$: dividendo pagado en el per\'iodo anterior.
  \item $c_i$: \textbf{factor de velocidad de ajuste (speed-of-adjustment factor)}, la
        fracci\'on de la brecha entre dividendo objetivo y dividendo previo que la empresa
        intenta reflejar en el pago corriente.
  \item $a_i$: constante que generalmente es positiva, capturando la mayor renuencia a
        \textit{reducir} que a \textit{aumentar} dividendos, y el deseo de crecimiento gradual.
  \item $u_{it}$: t\'ermino de error que absorbe discrepancias por preferencia de cifras
        redondeadas y otros factores no sistem\'aticos.
\end{itemize}

\begin{intuicionmat}
La ecuaci\'on dice: la empresa no salta directamente al dividendo ideal ($D^*_{it}$);
en cambio, mueve su dividendo en una fracci\'on $c_i$ de la distancia entre d\'onde est\'a
($D_{i,t-1}$) y d\'onde deber\'ia estar ($D^*_{it}$). Si $c_i = 1$, el ajuste es completo
en un a\~no; si $c_i = 0.25$, la empresa cierra s\'olo el 25\,\% de la brecha cada a\~no.
La constante $a_i > 0$ refleja que las empresas tienden a subir dividendos aunque sea un poco,
incluso cuando el ``ajuste puro'' indicar\'ia no cambiarlos.
\end{intuicionmat}

\subsection{Reformulaci\'on Econom\'etrica (Econometric Reformulation)}

Sustituyendo $D^*_{it} = r_i P_{it}$ en la ecuaci\'on (1) y reagrupando t\'erminos, se obtiene
la forma de regresi\'on directamente estimable:

\begin{equation}
  D_{it} \;=\; a_{it} + b\,P_{it} + d\,D_{i,t-1} + u_{it}
\end{equation}

con los par\'ametros compuestos:
\begin{align}
  b &= c \cdot r \\
  d &= (1 - c)
\end{align}

\begin{intuicionmat}
Esta forma dice que el dividendo de hoy es una combinaci\'on lineal de las ganancias
corrientes y el dividendo del per\'iodo anterior. El coeficiente $b$ recoge el efecto
conjunto de la raz\'on objetivo y la velocidad de ajuste; $d$ recoge la ``memoria'' del
sistema: cu\'anto persiste el dividendo pasado en el presente. Si $d$ es cercano a 1,
los dividendos son muy persistentes; si $d$ es cercano a 0, el ajuste es casi inmediato.
\end{intuicionmat}

\subsection{Ajuste del Modelo a los Datos del Campo}

El modelo describe con precisi\'on el comportamiento observado:

\begin{itemize}
  \item 26 de las 28 empresas ten\'ian un valor bien definido de $r_i$ (raz\'on objetivo).
  \item 20 empresas ten\'ian valores razonablemente definidos de $c_i$; las 6 restantes
        mostraban mayor flexibilidad.
  \item 22 empresas ajustaban dividendos a\~no a a\~no; las otras 4 lo hac\'ian cada 2--3 a\~nos.
  \item Aproximadamente el \textbf{85\,\% de los a\~nos-empresa} de las 26 empresas con
        est\'andares claros se explica con solo discrepancias moderadas mediante este modelo.
  \item El porcentaje sube a m\'as del 90\,\% si se permite una tolerancia de seis meses
        adicional y no se contabilizan las reducciones no realizadas.
\end{itemize}

\begin{intuicion}
El modelo con s\'olo dos par\'ametros por empresa ($r_i$ y $c_i$) captura cuatro quintas
partes del comportamiento observado. Esto es notable dado que las decisiones de dividendos
involucran factores cualitativos, personalidades y contextos \'unicos a cada empresa.
\end{intuicion}

%%% =========================================================
\section{Pruebas Estad\'isticas: Predicci\'on de Dividendos de Posguerra}
%%% =========================================================

\subsection{Estrategia de Validaci\'on}

Para evaluar la capacidad predictiva del modelo, Lintner procede as\'i:

\begin{enumerate}
  \item Ajusta la ecuaci\'on (2) a datos \textbf{de preguerra} del universo corporativo
        (datos de cuentas nacionales, 1918--1941, excluyendo 1936--1937 por el impuesto
        sobre utilidades no distribuidas).
  \item Usa los par\'ametros estimados para \textbf{pronosticar dividendos de posguerra}
        (1946--1954).
  \item Compara los errores de predicci\'on con los de cuatro modelos alternativos propuestos
        en la literatura y con un \textbf{modelo na\'ive (na\"ive model)} (que supone que el
        dividendo de este a\~no igual al del a\~no anterior).
\end{enumerate}

\begin{intuicion}
Ajustar el modelo \textit{fuera de muestra} es la prueba m\'as exigente: si los par\'ametros
estimados con datos de 1918--1941 predicen bien los datos de 1946--1954 (un per\'iodo
hist\'oricamente muy diferente), el modelo capta algo estructural y estable del comportamiento
corporativo, no solo una regularidad casual del per\'iodo de estimaci\'on.
\end{intuicion}

\subsection{Ecuaciones Estimadas con Datos de Preguerra}

Las ecuaciones ajustadas al per\'iodo 1918--1941 para el universo de corporaciones son:

\medskip
\noindent Con ajuste por valuaci\'on de inventarios (\textit{inventory valuation adjustment}):
\begin{equation}
  D \;=\; 352.3 + 0.15\,P_w + 0.70\,D_{-1}
\end{equation}

\noindent Sin ajuste por valuaci\'on de inventarios:
\begin{equation}
  D \;=\; 106.0 + 0.145\,P_n + 0.788\,D_{-1}
\end{equation}

donde $P_w$ son las ganancias con ajuste de inventarios, $P_n$ son las ganancias sin ajuste,
y $D_{-1}$ es el dividendo rezagado un per\'iodo (en miles de millones de d\'olares).

\begin{intuicionmat}
En la primera ecuaci\'on, $b = 0.15$ y $d = 0.70$. Esto implica:
$c = 1 - d = 0.30$ (se cierra el 30\,\% de la brecha por a\~no) y
$r = b/c = 0.15/0.30 = 0.50$ (raz\'on objetivo del 50\,\%).
Estos valores son \textit{exactamente} consistentes con los hallazgos del trabajo de campo:
velocidad de ajuste moderada y raz\'on objetivo del 50\,\%.
\end{intuicionmat}

\subsection{Resultados de las Predicciones Fuera de Muestra}

\begin{explicacion}
La Tabla~1 del art\'iculo original (p.~110) compara el error de predicci\'on de posguerra del
modelo de Lintner con el de siete especificaciones alternativas. Las variables $P_w$, $P_n$,
$NW$ denotan ganancias con ajuste de inventarios, sin ajuste y valor neto (\textit{net worth}),
respectivamente. $S$ representa el superavit.
\end{explicacion}

Los resultados clave son:

\begin{itemize}
  \item El modelo de Lintner con ajuste de inventarios produce un \textbf{error algebraico
        medio (algebraic mean error)} de $-163.7$ millones de d\'olares: una
        subestimaci\'on del 2.0\,\% de los dividendos reales y sobreestimaci\'on del 2.2\,\%
        de las ganancias retenidas.
  \item El \textbf{error absoluto medio (mean absolute error)} es del 6.4\,\% para
        dividendos y 5.6\,\% para ganancias retenidas.
  \item El \textbf{error cuadr\'atico medio (root-mean-square error)} es del 7.3\,\% para
        dividendos.
  \item Estos errores son de un tercio a cinco sextos \textit{menores} que los de los modelos
        alternativos basados en ganancias corrientes y rezagadas propuestos por Tinbergen
        y Modigliani, o con superavit, o basados en el valor neto.
  \item El modelo es m\'as del cuadrup\'edruple mejor que el modelo na\'ive en t\'erminos
        de error algebraico medio.
  \item Los estad\'isticos de von Neumann (1.98 y 2.04) caen pr\'oximos a su valor esperado
        de 2.25, confirmando la aleatoriedad de los residuos.
\end{itemize}

\begin{intuicion}
El modelo de dos par\'ametros estimado con datos de hace m\'as de diez a\~nos supera
consistentemente a modelos m\'as complejos estimados con m\'as variables. Esto sugiere
que los dos par\'ametros capturan \textit{la f\'isica del sistema}, no solo el ruido del
per\'iodo de estimaci\'on.
\end{intuicion}

\subsection{Robustez: Variables Adicionales No Significativas}

Lintner verifica que a\~nadir las siguientes variables \textbf{no} mejora el modelo
ni produce coeficientes estad\'isticamente significativos:

\begin{itemize}
  \item Gastos en planta y equipo (o desembolsos de capital netos de depreciaci\'on).
  \item Cambio en inventarios f\'isicos.
  \item Cambio corriente en ganancias o dividendos m\'aximos previos.
  \item Exceso de fondos internos sobre desembolsos de capital.
\end{itemize}

Este resultado refuerza la tesis del campo: la gerencia \textit{ya incorpora} las
consideraciones de inversi\'on y liquidez en sus est\'andares de largo plazo ($r$ y $c$);
no las trata como shocks adicionales a\~no a a\~no.

\begin{intuicion}
El modelo no falla por omitir la inversi\'on o la liquidez. Falla en incluirlas porque
estas variables no a\~naden informaci\'on una vez que se conocen ganancias y dividendo previo.
La empresa ya ``previ\'o'' estas restricciones al fijar su pol\'itica.
\end{intuicion}

%%% =========================================================
\section{Implicaciones Macroecon\'omicas: Ahorro Corporativo y Estabilizaci\'on}
%%% =========================================================

\subsection{Propensiones Marginales a Ahorrar (Marginal Propensities to Save)}

A partir de las ecuaciones estimadas se derivan las \textbf{propensiones marginales a ahorrar
corporativas (corporate marginal propensities to save, MPS)}:

\medskip
\noindent \textbf{MPS de corto plazo (short-run MPS):}

Dado que $\Delta D = c(rP - D_{-1}) + a$, ante un aumento permanente en $P$, el dividendo
sube solo en $b = cr$ en el primer a\~no. Por tanto, la fracci\'on del incremento de ganancias
que se retiene es:

\begin{equation}
  \text{MPS}_{\text{corto plazo}} \;=\; 1 - b \;=\; 1 - cr
\end{equation}

Con los valores emp\'iricos $c \approx 0.30$ y $r \approx 0.50$, se obtiene $b \approx 0.15$,
de modo que:
\[
  \text{MPS}_{\text{corto plazo}} \approx 0.85
\]

\begin{intuicionmat}
En el a\~no en que las ganancias suben \$1, la empresa retiene \$0.85 y paga solo \$0.15
en dividendos adicionales. Esto ocurre porque el ajuste al dividendo objetivo es gradual.
La alta MPS de corto plazo tiene un efecto \textit{estabilizador} sobre el ciclo econ\'omico:
ante una ca\'ida de ganancias, los dividendos tampoco caen bruscamente, sosteniendo el ingreso
disponible de los hogares accionistas.
\end{intuicionmat}

\noindent \textbf{MPS de largo plazo (long-run MPS):}

En el largo plazo (equilibrio est\'atico), $D = rP$ y la fracci\'on retenida es:

\begin{equation}
  \text{MPS}_{\text{largo plazo}} \;=\; 1 - r \;\approx\; 1 - 0.60 \;=\; 0.40
\end{equation}

donde $r \approx 0.60$ es la \textbf{raz\'on objetivo promedio ponderada (weighted average
target pay-out ratio)} del universo corporativo.

\begin{intuicionmat}
En el largo plazo, cuando los dividendos han alcanzado su nivel objetivo, la empresa retiene
solo el 40\,\% de sus ganancias. Esto contrasta con el 85\,\% del corto plazo: a medida que
pasan los a\~nos despu\'es de un aumento de ganancias, los dividendos van subiendo y la
retenci\'on va cayendo del 85\,\% al 40\,\%.
\end{intuicionmat}

\subsection{El Subsistema Dividendos--Ganancias--Ganancias Retenidas}

El sistema \textbf{dividendos--ganancias--ganancias retenidas (dividends--profits--retained
earnings subsystem)} es internamente \textbf{muy estable (highly stable)} aunque se encuentra
en \textbf{desequilibrio continuo (continuous disequilibrium)}: los dividendos nunca llegan
completamente al nivel objetivo porque las ganancias cambian antes de que se complete el ajuste.

\begin{intuicion}
El sistema es como un flotador en el agua: siempre trata de volver al equilibrio (raz\'on
objetivo), pero las olas (cambios en ganancias) lo sacan del equilibrio antes de que llegue.
El resultado es un movimiento continuo alrededor del equilibrio, nunca en \'el.
\end{intuicion}

\subsection{Pol\'itica de Posguerra: Los Dividendos no Fueron Deprimidos por Impuestos}

La evidencia es consistente con el juicio de que los dividendos de posguerra \textbf{no
fueron deprimidos} por debajo de las expectativas normales (en t\'erminos de ganancias
netas despu\'es de impuestos y pol\'iticas hist\'oricas) por la elevada carga tributaria
de la posguerra. Tampoco hay evidencia de que la raz\'on objetivo de largo plazo haya
cambiado entre la preguerra y la posguerra.

La explicaci\'on de los bajos \textbf{pay-out ratios} de posguerra no es una retenci\'on
discrecional para financiar la inversi\'on, sino que los gastos en inversi\'on f\'isica
han estado hist\'oricamente correlacionados con las ganancias corrientes y el flujo de
fondos internos, y las empresas han incorporado esa relaci\'on en sus pol\'iticas de dividendos.

\begin{intuicion}
Las empresas de posguerra no pagaron menos dividendos porque tuvieran m\'as proyectos de
inversi\'on: pagaron menos porque sus ganancias netas (ya despu\'es de impuestos) eran la
base del c\'alculo, y los impuestos redujeron esa base. La inversi\'on no fue una causa
independiente del bajo pay-out ratio.
\end{intuicion}

%%% =========================================================
\section{Resumen y Conclusiones}
%%% =========================================================

El art\'iculo de Lintner (1956) constituye la piedra fundacional del estudio moderno de la
pol\'itica de dividendos corporativos. Sus contribuciones centrales son:

\medskip
\noindent \textbf{1. Hallazgos del trabajo de campo:}
\begin{itemize}
  \item Los gerentes no optimizan el dividendo de cada per\'iodo de forma independiente;
        parten siempre de la tasa vigente y deciden si y cu\'anto cambiarla.
  \item Las ganancias netas corrientes son el \textit{\'unico} factor consistente y dominante
        en las decisiones de dividendos a lo largo de empresas y a\~nos.
  \item Las empresas tienen raz\'ones objetivo de pago ($r$) y velocidades de ajuste ($c$)
        estables en el tiempo, que resumen sus pol\'iticas de dividendo.
\end{itemize}

\noindent \textbf{2. El modelo te\'orico:}
\begin{equation*}
  \Delta D_{it} = a_i + c_i\!\left(r_i P_{it} - D_{i,t-1}\right) + u_{it}
\end{equation*}
Dos par\'ametros por empresa capturan el 85\,\% del comportamiento observado de dividendos.
La reformulaci\'on como $D_{it} = a + b P_{it} + d D_{i,t-1} + u_{it}$ es directamente
estimable por m\'inimos cuadrados ordinarios.

\noindent \textbf{3. Validaci\'on estad\'istica:}
\begin{itemize}
  \item Estimado con datos de preguerra (1918--1941), el modelo predice dividendos de
        posguerra (1946--1954) con errores de apenas 2\,\%, superando a todos los modelos
        alternativos de la literatura.
  \item Las variables de inversi\'on y liquidez no mejoran el ajuste: ya est\'an
        incorporadas en los par\'ametros de la pol\'itica de largo plazo.
\end{itemize}

\noindent \textbf{4. Implicaciones macroecon\'omicas:}
\begin{itemize}
  \item La \textbf{MPS de corto plazo} es $\approx 85\,\%$: los dividendos se ajustan lentamente
        a los cambios en ganancias, lo que da al subsistema un efecto estabilizador sobre el ciclo.
  \item La \textbf{MPS de largo plazo} es $\approx 40\,\%$: en equilibrio, las empresas
        distribuyen alrededor del 60\,\% de sus ganancias.
  \item Los altos impuestos de posguerra no deprimieron los dividendos de forma discrecional;
        simplemente redujeron la base (ganancias netas) sobre la que se aplica la raz\'on objetivo.
\end{itemize}

\noindent \textbf{5. Legado:}
El \textbf{modelo de Lintner (Lintner model)} es la referencia est\'andar para la econom\'ia
de la pol\'itica de dividendos. Su ecuaci\'on de ajuste parcial fue la base emp\'irica sobre
la que se construyeron posteriormente trabajos como el de Fama y Babiak (1968) y sigue siendo
el punto de partida en la literatura de finanzas corporativas sobre distribuci\'on de ingresos.

\end{document}
